% Options for packages loaded elsewhere
% Options for packages loaded elsewhere
\PassOptionsToPackage{unicode}{hyperref}
\PassOptionsToPackage{hyphens}{url}
\PassOptionsToPackage{dvipsnames,svgnames,x11names}{xcolor}
%
\documentclass[
  letterpaper,
  DIV=11,
  numbers=noendperiod]{scrartcl}
\usepackage{xcolor}
\usepackage[margin=2.5cm]{geometry}
\usepackage{amsmath,amssymb}
\setcounter{secnumdepth}{5}
\usepackage{iftex}
\ifPDFTeX
  \usepackage[T1]{fontenc}
  \usepackage[utf8]{inputenc}
  \usepackage{textcomp} % provide euro and other symbols
\else % if luatex or xetex
  \usepackage{unicode-math} % this also loads fontspec
  \defaultfontfeatures{Scale=MatchLowercase}
  \defaultfontfeatures[\rmfamily]{Ligatures=TeX,Scale=1}
\fi
\usepackage{lmodern}
\ifPDFTeX\else
  % xetex/luatex font selection
\fi
% Use upquote if available, for straight quotes in verbatim environments
\IfFileExists{upquote.sty}{\usepackage{upquote}}{}
\IfFileExists{microtype.sty}{% use microtype if available
  \usepackage[]{microtype}
  \UseMicrotypeSet[protrusion]{basicmath} % disable protrusion for tt fonts
}{}
\makeatletter
\@ifundefined{KOMAClassName}{% if non-KOMA class
  \IfFileExists{parskip.sty}{%
    \usepackage{parskip}
  }{% else
    \setlength{\parindent}{0pt}
    \setlength{\parskip}{6pt plus 2pt minus 1pt}}
}{% if KOMA class
  \KOMAoptions{parskip=half}}
\makeatother
% Make \paragraph and \subparagraph free-standing
\makeatletter
\ifx\paragraph\undefined\else
  \let\oldparagraph\paragraph
  \renewcommand{\paragraph}{
    \@ifstar
      \xxxParagraphStar
      \xxxParagraphNoStar
  }
  \newcommand{\xxxParagraphStar}[1]{\oldparagraph*{#1}\mbox{}}
  \newcommand{\xxxParagraphNoStar}[1]{\oldparagraph{#1}\mbox{}}
\fi
\ifx\subparagraph\undefined\else
  \let\oldsubparagraph\subparagraph
  \renewcommand{\subparagraph}{
    \@ifstar
      \xxxSubParagraphStar
      \xxxSubParagraphNoStar
  }
  \newcommand{\xxxSubParagraphStar}[1]{\oldsubparagraph*{#1}\mbox{}}
  \newcommand{\xxxSubParagraphNoStar}[1]{\oldsubparagraph{#1}\mbox{}}
\fi
\makeatother


\usepackage{longtable,booktabs,array}
\newcounter{none} % for unnumbered tables
\usepackage{calc} % for calculating minipage widths
% Correct order of tables after \paragraph or \subparagraph
\usepackage{etoolbox}
\makeatletter
\patchcmd\longtable{\par}{\if@noskipsec\mbox{}\fi\par}{}{}
\makeatother
% Allow footnotes in longtable head/foot
\IfFileExists{footnotehyper.sty}{\usepackage{footnotehyper}}{\usepackage{footnote}}
\makesavenoteenv{longtable}
\usepackage{graphicx}
\makeatletter
\newsavebox\pandoc@box
\newcommand*\pandocbounded[1]{% scales image to fit in text height/width
  \sbox\pandoc@box{#1}%
  \Gscale@div\@tempa{\textheight}{\dimexpr\ht\pandoc@box+\dp\pandoc@box\relax}%
  \Gscale@div\@tempb{\linewidth}{\wd\pandoc@box}%
  \ifdim\@tempb\p@<\@tempa\p@\let\@tempa\@tempb\fi% select the smaller of both
  \ifdim\@tempa\p@<\p@\scalebox{\@tempa}{\usebox\pandoc@box}%
  \else\usebox{\pandoc@box}%
  \fi%
}
% Set default figure placement to htbp
\def\fps@figure{htbp}
\makeatother





\setlength{\emergencystretch}{3em} % prevent overfull lines

\providecommand{\tightlist}{%
  \setlength{\itemsep}{0pt}\setlength{\parskip}{0pt}}



 


\KOMAoption{captions}{tableheading}
\makeatletter
\@ifpackageloaded{caption}{}{\usepackage{caption}}
\AtBeginDocument{%
\ifdefined\contentsname
  \renewcommand*\contentsname{Table of contents}
\else
  \newcommand\contentsname{Table of contents}
\fi
\ifdefined\listfigurename
  \renewcommand*\listfigurename{List of Figures}
\else
  \newcommand\listfigurename{List of Figures}
\fi
\ifdefined\listtablename
  \renewcommand*\listtablename{List of Tables}
\else
  \newcommand\listtablename{List of Tables}
\fi
\ifdefined\figurename
  \renewcommand*\figurename{Figure}
\else
  \newcommand\figurename{Figure}
\fi
\ifdefined\tablename
  \renewcommand*\tablename{Table}
\else
  \newcommand\tablename{Table}
\fi
}
\@ifpackageloaded{float}{}{\usepackage{float}}
\floatstyle{ruled}
\@ifundefined{c@chapter}{\newfloat{codelisting}{h}{lop}}{\newfloat{codelisting}{h}{lop}[chapter]}
\floatname{codelisting}{Listing}
\newcommand*\listoflistings{\listof{codelisting}{List of Listings}}
\makeatother
\makeatletter
\makeatother
\makeatletter
\@ifpackageloaded{caption}{}{\usepackage{caption}}
\@ifpackageloaded{subcaption}{}{\usepackage{subcaption}}
\makeatother
\usepackage{bookmark}
\IfFileExists{xurl.sty}{\usepackage{xurl}}{} % add URL line breaks if available
\urlstyle{same}
\hypersetup{
  pdftitle={Pengaruh Penetrasi Akses Internet terhadap Indeks Pembangunan Manusia (IPM): Pendekatan Probabilitas dan Analisis Korelasi},
  pdfauthor={Alfitra Mumtaz Hanafi (0403251058); Choirun Nisa (NIM\_Nisa); Natasha Umaiza (NIM\_Natasha)},
  colorlinks=true,
  linkcolor={blue},
  filecolor={Maroon},
  citecolor={Blue},
  urlcolor={Blue},
  pdfcreator={LaTeX via pandoc}}


\title{Pengaruh Penetrasi Akses Internet terhadap Indeks Pembangunan
Manusia (IPM): Pendekatan Probabilitas dan Analisis Korelasi}
\author{Alfitra Mumtaz Hanafi (0403251058) \and Choirun Nisa
(NIM\_Nisa) \and Natasha Umaiza (NIM\_Natasha)}
\date{2026-06-04}
\begin{document}
\maketitle

\renewcommand*\contentsname{Table of contents}
{
\setcounter{tocdepth}{3}
\tableofcontents
}

\section{BAB I: PENDAHULUAN}\label{bab-i-pendahuluan}

\subsection{Latar Belakang}\label{latar-belakang}

Dalam era modern, definisi infrastruktur pembangunan tidak lagi terbatas
pada pembangunan fisik, melainkan telah berekspansi ke infrastruktur
digital. Akses internet kini menjadi kebutuhan fundamental yang
bertindak sebagai katalisator perputaran ekonomi, pemerataan pendidikan,
dan akses informasi medis. Kesenjangan digital (\emph{digital divide})
sering kali menjadi akar masalah ketimpangan kesejahteraan antarwilayah.
Wilayah dengan penetrasi internet yang rendah cenderung mengalami
stagnasi dalam peningkatan kualitas Sumber Daya Manusia (SDM). Oleh
karena itu, penelitian ini bertujuan untuk mengevaluasi secara statistik
probabilitas dan signifikansi korelasi ketersediaan akses internet
terhadap Indeks Pembangunan Manusia (IPM) menggunakan bahasa pemrograman
R.

\subsection{Tujuan Analisis}\label{tujuan-analisis}

Analisis ini memiliki tujuan antara lain:

\begin{enumerate}
\def\labelenumi{\arabic{enumi}.}
\tightlist
\item
  Mengidentifikasi dan menangani \emph{missing value} serta
  \emph{outlier} pada dataset pembangunan wilayah, khususnya pada
  variabel akses internet dan IPM.
\item
  Melakukan statistika deskriptif untuk memetakan persebaran data
  digitalisasi dan kesejahteraan penduduk.
\item
  Membangun visualisasi data yang informatif untuk melihat tren sebaran
  antarvariabel secara grafis.
\item
  Menguji distribusi probabilitas (normalitas) pada data IPM wilayah.
\item
  Membuktikan secara statistik hipotesis korelasi antara penetrasi akses
  internet dengan skor IPM.
\end{enumerate}

\section{BAB II: METODOLOGI}\label{bab-ii-metodologi}

\subsection{Dataset}\label{dataset}

Penelitian ini menggunakan dataset
\texttt{pembangunan\_wilayah\_missing\_outlier\_2.csv} yang
merepresentasikan kondisi indikator pembangunan makro di berbagai
wilayah. Fokus pengujian dibatasi pada variabel independen
\texttt{akses\_internet} (persentase ketersediaan fasilitas internet)
dan variabel dependen \texttt{ipm} (Indeks Pembangunan Manusia). Dataset
ini secara inheren memiliki \emph{missing value} dan \emph{outlier} yang
disimulasikan untuk menguji ketahanan model pra-pemrosesan data sebelum
uji hipotesis dieksekusi.

\subsection{Tools yang Digunakan}\label{tools-yang-digunakan}

Seluruh tahapan pembersihan dan pengujian data dieksekusi menggunakan
bahasa pemrograman \textbf{R}. Dokumen ini disusun menggunakan
\textbf{Quarto}, sebuah sistem penerbitan saintifik yang memungkinkan
integrasi narasi teks dan eksekusi kode (\emph{reproducible research})
secara instan ke dalam format PDF. Pustaka tambahan seperti
\texttt{ggplot2} diimplementasikan untuk membangun visualisasi grafis
yang komprehensif.

\subsection{Metode Analisis}\label{metode-analisis}

Pendekatan statistik yang diterapkan meliputi:

\begin{itemize}
\tightlist
\item
  \textbf{Penanganan Data:} Imputasi data kosong diselesaikan
  menggunakan \emph{Median Imputation} untuk meminimalisasi efek
  distorsi akibat nilai ekstrem, sementara \emph{outlier} dideteksi
  melalui batas \emph{Interquartile Range} (IQR).
\item
  \textbf{Uji Distribusi Probabilitas:} Pemenuhan asumsi normalitas data
  diuji melalui \emph{Shapiro-Wilk Test}.
\item
  \textbf{Uji Korelasi Hipotesis:} Pengujian signifikansi menggunakan
  \emph{Pearson Correlation Test} guna mengukur kekuatan dan arah
  hubungan antara variabel bebas dan terikat.
\end{itemize}

\section{BAB III: HASIL DAN
PEMBAHASAN}\label{bab-iii-hasil-dan-pembahasan}

\subsection{Pemahaman dan Pembersihan
Dataset}\label{pemahaman-dan-pembersihan-dataset}

Tahap pertama adalah memuat dataset ke dalam memori lingkungan kerja dan
memindai anomali data.

```\{r\} \#\textbar{} label: load-and-clean \#\textbar{} warning: false
\#\textbar{} message: false

\section{Membaca dataset}\label{membaca-dataset}

data \textless-
read.csv(``../dataset/pembangunan\_wilayah\_missing\_outlier\_2.csv'')

\section{Pengecekan Missing Value
Awal}\label{pengecekan-missing-value-awal}

na\_internet\_awal \textless- sum(is.na(data\(akses_internet))
na_ipm_awal <- sum(is.na(data\)ipm))

cat(``Jumlah NA Akses Internet sebelum diolah:'', na\_internet\_awal,
``\n'') cat(``Jumlah NA IPM sebelum diolah:'', na\_ipm\_awal, ``\n'')

\section{Penanganan Missing Value menggunakan Median
Imputation}\label{penanganan-missing-value-menggunakan-median-imputation}

data\(akses_internet[is.na(data\)akses\_internet){]} \textless-
median(data\(akses_internet, na.rm = TRUE)
data\)ipm{[}is.na(data\(ipm)] <- median(data\)ipm, na.rm = TRUE)




\end{document}
